\section{Protobuf Descriptor Design}\label{sec:proto-desc}

The intermediate encoding format used a custom extensionality theorem over a
descriptor structured as a list to show that descriptors were equal. This was
required since Lean's finite maps (unlike Rocq's \texttt{gmap}) can't be placed
inside of inductive types. That's a clunky choice which was not what I wanted
when I told Aristotle ``I want extensionality on my descriptors'', but I kinda
just lived with it. Well, the full protobuf descriptors we're taking a more
principled approach that can fall back on the vast support of Mathlib.

\subsection{Protobuf Descriptors in Lean}

Instead, the full protobuf descriptors take a different approach, using a
dependently typed list internally but exposing a view of the descriptor using a
proper finite map.

\begin{minted}{lean}
inductive ScalarType where
  | double | float | int32 | int64 | uint32 | uint64
  | sint32 | sint64 | fixed32 | fixed64 | sfixed32 | sfixed64
  | bool | string | bytes

inductive Cardinality where
  | singular | optional | repeated | oneof (group : Nat)

mutual
inductive Desc      where | mk (es : List ((_ : Int) × Field))
inductive Field     where | mk (card : Cardinality) (ty : FieldType)
inductive FieldType where | scalar (s : ScalarType) | msg (d : Desc)
end
\end{minted}

Outside of Lean, we'll use this syntax for the Protobuf descriptors.

\begin{figure}[H]
  \centering
  \begin{syntax}
    \categoryFromSet[Naturals]{n}{\mathbb{N}}

    \category[Scalar Type]{\tau_s}
    \alternative{\double}
    \alternative{\float}
    \alternative{\ints}
    \alternative{\intl}
    \groupleft{
        \alternative{\uints}
        \alternative{\uintl}
        \alternative{\sints}
        \alternative{\sintl}
    }
    \groupleft{
      \alternative{\fixeds}
      \alternative{\fixedl}
      \alternative{\sfixeds}
      \alternative{\sfixedl}
    }
    \groupleft{
      \alternative{\bool}
      \alternative{\str}
      \alternative{\byt}
    }

    \category[Cardinality]{c}
    \alternative{\sig}
    \alternative{\opt}
    \alternative{\rep}
    \alternative{\oneof{n}}

    \category[Field]{f}
    \alternative{\langle c, \tau_s \rangle}
    \alternative{\langle c, d \rangle}

    \category[Field List]{l}
    \alternative{\texttt{[]}}
    \alternative{(n, f) :: l}

    \category[Reserved Set]{r}
    \alternative{\varnothing_r}
    \alternative{n, r}

    \category[Descriptor]{d}
    \alternative{\llangle r, l \rrangle}
  \end{syntax}
  \caption{Pollux syntax for Protobuf descriptors.}\label{fig:proto-desc-syn}
\end{figure}

An alternative representation of the reserved set is to insert an artificial
empty type, which all other types can be updated to but cannot be updated to any
other type itself. Then ``removing'' a field is updating a field to this empty
type $\bot$ and that would prevent that field from ever being added again,
enforcing the same mechanism as the reserved field set.

The use of the sigma type is important, since it let's us tie into the existing
Mathlib infrastructure for building a finite map from from the list
internals. The basic idea is that the descriptor itself is \emph{sealed}, and we
can explode one layer at a time to get a view of the descriptor through the lens
of a finite map which should enable reasoning about the descriptor using the
theory from Mathlib. This approach also is designed to respect these
requirements:

\begin{itemize}
\item \textbf{Kernel Acceptance.} Clearly we need a representation that Lean can
  handle. Under this choice, we're not storing the map directly but are still
  able to indirectly access other theorems about it.
\item \textbf{Extensionality.} The existing formulation certain variants of the
  compatibility theorems state that we recover not an equivalent value, but an
  \emph{equal} value, which \texttt{Finmap} provides.
\item \textbf{Termination Measures.} Since we have nested messages, we need to
  ensure that we can safely traverse the descriptor and recurse on the
  descriptor. Pervious formats, like the intermediate one, defined a custom
  termination metric by noting that the depth nested descriptor must be smaller
  than the containing descriptor. This works well for tree-like descriptors, but
  Protobuf does support recursive and mutually-recursive messages definitions,
  which cannot be represented by this descriptor type.
\item \textbf{Forward Compatibility.} Protobuf can represent it's own
  descriptors, but not without recursive messages. Eventually we'll have to
  support getting a descriptor into the Lean formalization, and one option is to
  model a descriptor as a value of \texttt{descriptor.proto} and then build the
  descriptor from that information. This is common in real-world protobuf
  implementations, and the sealed design should still work for this approach.
\end{itemize}

There is a well-formedness requirement on descriptors, one coming from the
internal structure of the descriptor in Lean (Definition~\ref{def:dwf}) and the
other from the Protobuf compiler (Definition~\ref{def:dlegal}).

\vspace{1em}
\begin{definition}[Descriptor Well-Formedness]\label{def:dwf}
  A descriptor $d$ is well-formed, denoted $d \checkmark_w$, if all fields in
  the descriptor have a unique field number, the fields are stored in sorted
  order, and all nested descriptors are also well-formed.
\end{definition}

\vspace{1em}
\begin{definition}[Legal Descriptor]\label{def:dlegal}
  A descriptor $d$ is legal, denoted $d \checkmark_l$, if
  \begin{itemize}
  \item All field numbers are between 1 and $2^{29}-1$, excluding the range
    19,000 to 19,999, which are reserved by the Protobuf compiler.
  \item Message fields are always explicitly marked as optional. This conforms
    with Protobuf, which will generate a method to check message presence, and
    will allow us to initialize recursive messages.
  \item All nested descriptors within $d$ are also legal.
  \end{itemize}
\end{definition}

In preparation for needing to support recursive messages, Pollux uses
well-founded recursion on the size of the descriptor through the \texttt{get?}
interface. We define a function to compute the size of the descriptor below.

\vspace{1em}
\begin{definition}[Descriptor Size]\label{def:dsize}
  The size of a descriptor $d$ is defined as $\mathrm{size}_d$ below.
  \begin{mathpar}

    \mathrm{size}_d\left(\llangle r, l \rrangle\right) = 1 + \mathrm{size}_d(l)

    \mathrm{size}_d(\langle \_, \tau_s \rangle) = 1
    
    \mathrm{size}_d(\langle \_, d \rangle) = 1 + \mathrm{size}_d(d) 

    \mathrm{size}_d(l) =
    \begin{cases}
      0 & \text{if } l = \texttt{[]} \\
      \mathrm{size}_d(f) + \mathrm{size}_d(l') & \text{if } l = (\_, f) :: l'
    \end{cases}
  \end{mathpar}
\end{definition}

Finally, the \texttt{oneof} cardinality needs some explanation. When looking at
the syntactic structure of a Protobuf message in a \texttt{.proto} file, the
logical way to model a \texttt{oneof} is likely to have a set of fields or some
other similar grouped structure. However, this isn't how Protobuf actually
models a \texttt{oneof} internally, and there are good reasons for Pollux to
follow their lead. 

\begin{itemize}
\item Members of a \texttt{oneof} share a field number namespace with the rest
  of the message. You can't define a field with identifier $n$ in two
  \texttt{oneof}'s in the same message, and using the \texttt{oneof} cardinality
  like we have it lets the well-formedness condition handle this automatically.
\item A \texttt{oneof} field doesn't have a unique field number, so there is no
  key available to insert a dedicated \texttt{oneof} representation into the
  finite map the descriptor is modeled as.
\item A field in a \texttt{oneof} can't be repeated, and \texttt{oneof}'s can't
  be empty. Both of those requirements are structurally impossible.
\item The Protobuf underlying \texttt{oneof} index isn't evolution stable, since
  introducing a new \texttt{oneof} will shift the assigned index of everything
  later. This complicates the \texttt{oneof} compatibility rules, but at least
  by adopting the same representation, everything is transparent.
\item If a program called \texttt{get?} with a field in a \texttt{oneof}, the
  \texttt{get?} operation would need to explore the inside of the \texttt{oneof}
  (via a mechanism perhaps similar to \texttt{explode}) anyways.
\end{itemize}

Fortunately, checking equality on the natural number included with each
\texttt{oneof} cardinality provides the required information. Even with a nested
or grouped representation, the parser would still have to manually unset the
previously set field, so this representation doesn't add any unnecessary work.

\subsection{Sealed \& Unsealed}

While the descriptor internals use a sorted listed, the public interface should
respect a \emph{seal} by convention (lean doesn't have great mechanisms to mark
things private or opaque). The public interface is this: 

\begin{minted}{lean}
def explode (d : Desc) : Finmap (fun _ : Int => Field) :=
    d.entries.toFinmap
def get?    (d : Desc) (k : Int) : Option Field := d.explode.lookup k
\end{minted}

Where \texttt{explode} unwraps the outer-most descriptor into a Mathlib
\texttt{Finmap}, enabling proofs to use the established \texttt{Finmap}
theory. Because \texttt{Finmap} quotients over the ordering of values in the
map, the sorted-list invariant is invisible most consumers of the interface, as
long as they use the provided mechanisms when inserting into a descriptor. Some
operations do still require the well-formed predicate, such as extensional
equality, deletions and a helpful lemma which shows that if two descriptors have
the same exploded result, they must be the same descriptor.

Descriptors should be sealed. Looking at how the code for the intermediate
format is structured, descriptors are never consumed more than one layer at a
time. The serializer iterates values, not the descriptor. While looking at a
message, we check against fields as we encounter them in the value.

However, values do need to iterated over and are the basis for round-trip
proofs. Sealing the values within the parser and serializer implementation would
create more work than we desire, although the theorems will still be able to use
a variant of the explode operator if it makes more sense.

\subsection{Considerations about Ordering \& Determinism}

The Protobuf spec permits serializing and parsing fields in any order, which
seems at odd with the choice to represent the fields in a sorted list. It should
be possible to explode a layer and use the undefined iteration order of the
\texttt{Finmap} to be as generalize as possible with regards to that Protobuf
spec. However,

\begin{itemize}
\item Having the serializer return data in any order would mean that it isn't a
  pure function, which is difficult to represent in Lean. We could quotient over
  the ordering of the fields, but then we'd have to prove that every other
  function consuming the serialized output also respects the quotient type,
  which is possible but a lot of work.
\item The exact output from the parser is descriptor dependent as well as value
  dependent (recall that we can't know what value is encoded without needing to
  to a descriptor, see Section~\ref{sec:philosophy}). Having both align as an
  ordered list reduces map equivalences to syntactic list equality.
\item The Protobuf spec is also adding a requirement to the parser, that is has
  to be able to handle encodings which were produced with the fields in an
  arbitrary order, with last field wins semantics. The existing \texttt{ParseOk}
  definition used for the intermediate formalization only requires the parser to
  be correct on the output of our own serializer, not any possible encoding. It
  will need to be extended with a relational specification which allows the
  maximum possible amount of non-determinism. The soundness requirement for this
  specification would require that our serializer produces values that satisfy
  the relational spec while completeness would require that any encoding in the
  relational specification can be parsed correctly.
\end{itemize}

\subsection{New Protobuf Features}

There are three new, notable features in Protobuf that were missing from the
intermediate formate entirely.

\begin{itemize}
\item \textbf{Cardinality.} The intermediate format only had scalar values, but
  full Protobuf has repeated values. (Optional values and implicit values are
  already extremely similar, so the addition of repeated values is the most
  important). 
\item \textbf{Maps.} Maps are easy to support since we can just desugar them, as
  described in Section~\ref{sec:proto}, with a few additions for the
  restrictions on map key type and the like.
\item \textbf{Oneof.} On the wire, a \texttt{oneof} is invisible. There is no
  information included in the wire level format to describe that a field is in a
  \texttt{oneof}, it's another one of those things that we have to have a
  descriptor to know about. Individual fields still have their own field number,
  which is encoded, that comes from the name space of the whole message. The
  \texttt{oneof} really adds a parser semantic requirement about clearing the
  old field if we encounter a new encoded field in the same \texttt{oneof}, and
  overall constraint on the final message that at most one of the fields will be
  set. The \texttt{oneof} is then represented as a cardinality, with some extra
  information to determine which \texttt{oneof} in the message this field
  belongs to. This exactly mirrors how the \texttt{descriptor.proto} operates,
  and we can use natural number equality to scan the value and clear old set
  fields when needed. Since each field has exactly one cardinality, this already
  enforces that we can't have repeated fields inside a oneof and that a oneof
  can't be empty.
\end{itemize}

\subsection{Omitted Protobuf Features}

At the moment, these features are \emph{not} modeled in Pollux. Adding support
for them should be relatively straight-forward. If I've already developed
relevant rules for them, they might be included in a different color.

\textbf{Some of these features will likely be added before the work is complete.}

\begin{itemize}
\item \textbf{Reserved field numbers.} Protobuf doesn't recommend reusing field
  numbers for new fields (i.e.\ removing a field, then later adding a field with
  that same number), and that also poses a challenge to Pollux. If a field were
  added which uses and old field number, Pollux would assume that the intent is
  for the new field to be compatible with the old one. Since we expect the
  compatibility relation to be transitive, any evolving derivation of that
  message would attempt to force a connection between two fields that are
  semantically unrelated.

  Fortunately, Protobuf has a mechanism for this, reserved field numbers. You
  can set specific field numbers as reserved, then the protobuf compiler will
  ensure that you don't add a field with that number to the message in the
  future. Pollux will also enforce this, even though the Lean formalization
  doesn't track it \emph{yet}.

\item \textbf{Field names.} Every Protobuf field has a name, which is used in
  the generated code from \texttt{protoc}. These names aren't recorded at all in
  the wire format, but do impact the compatibility of a change with the rest of
  your code. At the moment, that's considered beyond the scope of Pollux.

\item \textbf{Reserved field names.} Like reserved field numbers, you can also
  reserve field names. Pollux's stance is that this is outside the domain we're
  targeting, but could be added with mechanisms similar to the reserved field
  number.

\item \textbf{Recursive Messages.} Protobuf messages (and thus descriptors) can
  be recursive. The current lean formalization enforces a tree structure on the
  descriptors. This is nice for recursion over the descriptors, but does mean we
  can't model all descriptors. The Protobuf data set does contain recursive and
  mutually recursive descriptors. It seems that these are more common in tools
  which analyze Protobuf descriptors, likely because \texttt{descriptor.proto},
  which models Protobuf descriptors as a Protobuf struct, is recursive. If we
  want to transmit protobuf descriptors into the Pollux lean formalization using
  \texttt{descriptor.proto} (a reasonable choice), we'll have to expand the
  descriptor representation to be a flat-map of named messages. See
  Table~\ref{tab:recur-summary} and Table~\ref{tab:recur-repos} for information
  about recursive messages in the Protobuf data set.

\item \textbf{Enums.} Protobuf supports enums, which shouldn't be hard to
  support in the future, but aren't present currently.

\item \textbf{Groups.} This is a Protobuf 2 feature, which still has some
  reserved encoding wire type implications. It's simply not used in Protobuf 3.

\item \textbf{Explicit Maps.} Maps can be desugared to a \texttt{repeated}
  key-value pair (See Section~\ref{sec:proto}), so no explicit representation is
  needed. Ok, this is technically \emph{not} true, since maps have additional
  requirements about the type of the key, but if we were to desugar a map to
  it's internal representation, the user would need to also desugar it to create
  something which might be problematic.

\item \textbf{Packed Scalars.} Not really a concern of the descriptor, although
  there are serious consequences of this.
\end{itemize}

\begin{table}[H]
\centering
\begin{tabular}{lcc}
    \toprule
    Metric & Count & Percentage \\
    \midrule
    Messages & 57,461 & \\
    Recursive & 831 & 1.45\% \\
    \bullet{} self-recursive & 425 & 0.74\% \\
    \bullet{} mutually-recursive & 447 & 0.78\% \\
    Contains recursion & 3,893 & 6.78\% \\
    \bottomrule
\end{tabular}
\caption[Summary of recursive protobuf messages.]{Summary statistics of
    recursive descriptor usage in the Protobuf data set. Here ``recursive''
    means directly self-recursive or in a mutually recursive cycle while
    ``contains recursion'' is any message which includes a recursive message in
    it's transitive closure. Any message which contains recursive cannot be
    represented by the current descriptor model. There are 533
    mutually-recursive cycles with the longest containing 22 messages. Note
    that there are 41 messages which are both self-recursive and
    mutually-recursive.}\label{tab:recur-summary}
\end{table}

\begin{table}[H]
\centering
\begin{tabular}{lccccc}
    \toprule
    Repository & Messages & Recursive & \% & Contains & \% \\
    \midrule
    \texttt{protocolbuffers/protobuf-go} & 605 & 110 & 18.18\% & 151 & 24.96\% \\
    \texttt{cloudwego/kitex} & 34 & 4 & 11.76\% & 7 & 20.59\% \\
    \texttt{go-kratos/kratos} & 94 & 8 & 8.51\% & 17 & 18.09\% \\
    \texttt{bufbuild/buf} & 2,922 & 210 & 7.19\% & 371 & 12.70\% \\
    \texttt{google/protobuf.dart} & 503 & 30 & 5.96\% & 59 & 11.73\% \\
    \texttt{protocolbuffers/protobuf} & 2,372 & 172 & 7.25\% & 242 & 10.20\% \\
    \texttt{googleapis/googleapis} & 49,370 & 289 & 0.59\% & 3,022 & 6.12\% \\
    \texttt{streamdal/plumber} & 286 & 2 & 0.70\% & 13 & 4.55\% \\
    \texttt{google/rejoiner} & 50 & 0 & 0.00\% & 2 & 4.00\% \\
    \texttt{googleapis/google-cloud-go} & 134 & 3 & 2.24\% & 3 & 2.24\% \\
    \texttt{uber/prototool} & 469 & 3 & 0.64\% & 6 & 1.28\% \\
    \texttt{asynkron/protactor-go} & 149 & 0 & 0.00\% & 0 & 0.00\% \\
    \texttt{colinmarc/hdfs} & 443 & 0 & 0.00\% & 0 & 0.00\% \\
    \texttt{twitchtv/twirp} & 30 & 0 & 0.00\% & 0 & 0.00\% \\
    \bottomrule
\end{tabular}
\caption[Per-repo recursive message breakdown]{A pre-repo breakdown of
    recursive message and messages which contain recursion. Notice that tools
    which work with protobuf descriptors have a higher percentage of recursive
    messages than project which just use Protobuf. The aggregate numbers in
    Table~\ref{tab:recur-summary} are dominated by
    \texttt{googleapis/googleapis} since it contains 86\% of the data set. I
    don't know what's happening with \texttt{google/rejoiner}, but it's
    possible some messages didn't compile correctly.}\label{tab:recur-repos}
\end{table}

% Local Variables:
% citar-bibliography: ("../pollux.bib")
% TeX-master: "../pollux.tex"
% TeX-command-default: "Make"
% jinx-local-words: "protobuf"
% End:
